\section{Recall of smooth affine algebraic geometry over $\C$}
Let $R$ be an integral commutative $\C$-algebra of finite type. Set
\[
  X := \Spec R,
\]
 the affine algebraic variety associated to $R$, endowed with the Zariski topology.
For an ideal $I\subseteq R$, write
\[
  V(I) := \{\, \mathfrak p\in \Spec R \mid \mathfrak p\supseteq I \,\},
\]
so that a subset $V\subseteq X$ is closed iff $V=V(I)$ for some $I$.

We denote the module of K"ahler differentials of $R$ over $\C$ by
\[
  \Omega_R = \Omega_X = \Omega_{R/\C}.
\]
\begin{definition}[K\"ahler differentials]\label{def:kahler}
The $R$-module $\Omega_{R/\C}$ is generated by formal symbols $\mathrm df$ with $f\in R$,
subject to the relations, for all $f,g\in R$ and $a,b\in\C$,
\begin{align*}
  \mathrm d(fg) &= f\,\mathrm dg + g\,\mathrm df \qquad\text{(Leibniz rule)},\\
  \mathrm d(af+bg) &= a\,\mathrm df + b\,\mathrm dg.
\end{align*}
\end{definition}
\begin{exercise}
Show that $\mathrm dc=0$ for every constant $c\in\C$.
\end{exercise}
\begin{proof}[Solution to the exercise]
Using linearity and the Leibniz rule,
\(\mathrm d(1)=\mathrm d(1\cdot 1)=1\,\mathrm d1+1\,\mathrm d1=2\,\mathrm d1\), hence
\(\mathrm d1=0\). For $c\in\C$ we have
\(\mathrm dc=\mathrm d(c\cdot 1)=c\,\mathrm d1+1\,\mathrm dc=\mathrm dc\), and the only possibility is
\(\mathrm dc=0\).
\end{proof}

Define the $R$-module
\[
  \T_X := \Hom_R\!\big(\Omega_{R/\C},\, R\big) \cong \Der_{\C}(R,R),
\]
called the module of $\C$-derivations on $X$ (the vector fields on $X$).

\begin{definition}[Smoothness]
We say that $R$ (or $X=\Spec R$) is \emph{smooth} if $\Omega_{R/\C}$ is a free $R$-module of
finite rank. In this case, the rank
\[
  n := \operatorname{rank}_R(\Omega_{R/\C})
\]
which is equal to the topological dimension of $X$.

\begin{proof}[Prove that $\mathrm{rank}_R(\Omega_{R/\mathbb{C}}) = \dim X$]
Let $K = \mathrm{Frac}(R)$ denote the fraction field of $R$. Since $R$ is integral.

First, recall a fundamental result from field theory \cite[Chapter 2, Proposition 8.6A.]{hartshorne1977algebraic}: for a field extension $K/\mathbb{C}$, the module of Kähler differentials $\Omega_{K/\mathbb{C}}$ is a vector space over $K$, and its dimension is equal to the transcendence degree of the extension:
\[
    \dim_K (\Omega_{K/\mathbb{C}}) = \mathrm{tr.deg}_{\mathbb{C}}(K).
\]

One should note that $\mathbb{C}$ is of characteristic 0, which ensures all extensions are separable, even when extensions are not algebraic.

By the definition of the dimension of an algebraic variety \cite[Chapter 1, Proposition 1.8A.(a)]{hartshorne1977algebraic}, we have:
\[
    \dim X = \mathrm{tr.deg}_{\mathbb{C}}(K).
\]

Next, we relate the differentials on $R$ to those on $K$ via localization. We have the canonical isomorphism Chapter 2, Proposition 8.2A. \cite[Chapter 2, Proposition 8.2A.]{hartshorne1977algebraic}:
\[
    \Omega_{K/\mathbb{C}} \cong \Omega_{R/\mathbb{C}} \otimes_R K.
\]
The hypothesis states that $\Omega_{R/\mathbb{C}}$ is a \textbf{free $R$-module} of rank $n$. A key property of free modules is that their rank is preserved under base change to the fraction field. Therefore:
\[
    \dim_K (\Omega_{R/\mathbb{C}} \otimes_R K) = \mathrm{rank}_R (\Omega_{R/\mathbb{C}}) = n.
\]

Combining these equalities, we conclude:
\[
    n = \dim_K (\Omega_{R/\mathbb{C}} \otimes_R K) = \dim_K (\Omega_{K/\mathbb{C}}) = \mathrm{tr.deg}_{\mathbb{C}}(K) = \dim X.
\]  

\end{proof}

\end{definition}

If $X$ is smooth, a tuple $(x_1,\dots,x_n)$ with $x_i\in R$ is called an \emph{algebraic
coordinate system} if
\[
  \Omega_{R/\C} \cong \bigoplus_{i=1}^n R\,\mathrm d x_i.
\]
Equivalently, the dual basis gives
\[
  \T_X \cong \bigoplus_{i=1}^n R\,\partial_{x_i},
\]
where $\partial_{x_i}$ are the vector fields dual to $\mathrm d x_i$.

\begin{example}[Polynomial ring]
Let $R=\C[x_1,\dots,x_n]$. Then
\[
  \Omega_{R/\C} \cong \bigoplus_{i=1}^n R\,\mathrm d x_i,
\]
so $X=\Spec R\simeq \A^n_{\C}$ is smooth and $(x_1,\dots,x_n)$ is an algebraic coordinate system.
\end{example}

From now on assume $X=\Spec R$ is smooth affine and $(x_1,\dots,x_n)$ are algebraic coordinates.
The symmetric algebra of vector fields is
\[
  \Sym^{\bullet}\T_X = \bigoplus_{m\ge 0}\Sym^{m}\T_X \simeq R[\xi_1,\dots,\xi_n],
\]
canonically an $R$-algebra (identify $\xi_i$ with the class of $\partial_{x_i}$). Define the
cotangent bundle
\[
  T^{\ast}X := \Spec\, \Sym^{\bullet}\T_X \xrightarrow{\quad\pi\quad} X.
\]
and $\pi$ is induced by the $R$-algebra map $R = \Sym^{0}\T_X \subset \bigoplus_{m\ge 0}\Sym^{m}\T_X$.


\section{D-modules on smooth affine varieties}
Let $\End_{\C}(R)$ denote the $\C$-algebra of all $\C$-linear maps $R\to R$. There are natural
embeddings
\[
  R\hookrightarrow \End_{\C}(R),\quad f\longmapsto (r\mapsto fr),
\]
and
\[
  \T_X\hookrightarrow \End_{\C}(R),\quad v\longmapsto v':=[f\mapsto v(\mathrm d f)],
  \begin{tikzcd}[row sep=3em, column sep=1.5em]
      R \arrow[rr, "d"] \arrow[dr, "v'"]  % v' 在箭头上方（默认）
      & & \Omega_R \arrow[dl, "v"]       % v 在箭头下方（不使用 ' 换边）
      \\
      & R &
  \end{tikzcd}.
\]

\begin{definition}[Ring of differential operators]
The ring $\D_X$ of algebraic differential operators on $X$ is the $\C$-subalgebra of
$\End_{\C}(R)$ generated by $R$ and $\T_X$.
\end{definition}

\begin{example}[Weyl algebra]
If $X=\A^n_{\C}=\Spec \C[x_1,\dots,x_n]$, then
\[
  \D_X \cong \C[x_1,\dots,x_n]\langle \partial_{x_1},\dots,\partial_{x_n}\rangle,
\]
with relations
\[
  [\partial_{x_i},x_j]=\delta_{ij},\quad [x_i,x_j]=0,\quad [\partial_{x_i},\partial_{x_j}]=0.
\]
\end{example}

Thus $\D_X$ is non-commutative. Every operator $P\in\D_X$ can be written in multi-index notation as
\[
  P=\sum_{\underline{\alpha}\in\mathbb N^n} f_{\underline{\alpha}}\,\partial^{\underline{\alpha}},\qquad
  \partial^{\underline{\alpha}}:=\partial_{x_1}^{\alpha_1}\cdots\partial_{x_n}^{\alpha_n},\quad f_{\underline{\alpha}}\in R.
\]

\paragraph{Basic modules.}\mbox{}\par

\noindent 1.  The coordinate ring $R$ is a left $\D_X$-module where$\partial_{x_i}\cdot r = \partial_{x_i}(\mathrm d r)\in R$.

\noindent Equivalently, writing $\mathrm dr=\sum_i (\partial_{x_i}r)\,\mathrm d x_i$, we have $\partial_{x_i}r$ the
$i$-th coefficient.\\

\par\noindent 2.  The top exterior power
\[
  \omega_X := \bigwedge{}^{n}\Omega_{R/\C} \cong R\,\mathrm d x_1\wedge\cdots\wedge\mathrm d x_n
\]
is the \textbf{dualizing module} of $R$. It carries a natural right $\D_X$-module structure via the transpose
anti-involution $t\colon \D_X\to\D_X$ determined by
\[
  \prescript{t}{}{x_i}=x_i,\qquad \prescript{t}{}{\partial_{x_i}}=-\partial_{x_i}.
\]
If $P=\sum_{\underline{\alpha}} f_{\underline{\alpha}}\,\partial^{\underline{\alpha}}$, then
\[
  \prescript{t}{}{P}=\sum_{\underline{\alpha}} (-\partial)^{\underline{\alpha}} f_{\underline{\alpha}},
\]
and the action is
\[
  (r\,\mathrm d x_1\wedge\cdots\wedge\mathrm d x_n)\cdot P := \left(\prescript{t}{}{P}\cdot r\right)\,\mathrm d x_1\wedge\cdots\wedge\mathrm d x_n.
\]

\subsection*{Side change and opposite ring}
Let $A$ be a ring and let $A^{\mathrm{op}}$ denote the opposite ring. i.e. $A^{\op} = \{ a^{\op} \mid a \in A \}$, $a^{\op} \cdot b^{\op} = (ba)^{\op}$. Then right $A$-modules are the
same as left $A^{\mathrm{op}}$-modules.

\begin{lemma}\label{lem:DXop}
There is a natural isomorphism of rings
\[
  \D_X^{\mathrm{op}} \simeq \omega_X \otimes_R \D_X \otimes_R \omega_X^{\!*},\qquad
\text{where } \omega_X^{\!*}:=\Hom_R(\omega_X,R) \cong R\,(\mathrm d x_1\wedge\cdots\wedge\mathrm d x_n)^{-1}.
\]
\end{lemma}
\begin{proof}
Define
\[
  \Phi(P^{\mathrm{op}}):=\mathrm d x_1\wedge\cdots\wedge\mathrm d x_n\;\otimes\; \prescript{t}{}{P}\;\otimes\;
  (\mathrm d x_1\wedge\cdots\wedge\mathrm d x_n)^{-1}.
\]
Then for $P,Q\in\D_X$ we have
\begin{align*}
  \Phi(P^{\mathrm{op}}Q^{\mathrm{op}})
  &= \Phi((QP)^{\mathrm{op}})\\
  &= \mathrm d x\;\otimes \prescript{t}{}{(QP)}\otimes (\mathrm d x)^{-1}\\
  &= \mathrm d x\;\otimes \prescript{t}{}{P}\prescript{t}{}{Q}\otimes (\mathrm d x)^{-1}\\
  &=: \left(\mathrm d x\;\otimes \prescript{t}{}{P}\otimes (\mathrm d x)^{-1}\right)
    \cdot \left(\mathrm d x\;\otimes \prescript{t}{}{Q}\otimes (\mathrm d x)^{-1}\right)\\
  &= \Phi(P^{\mathrm{op}})\,\Phi(Q^{\mathrm{op}}),
\end{align*}

so $\Phi$ is a ring homomorphism. 
\end{proof}

\begin{remark}
  $\D_X$ is a $(R,R)$-bimodule in the natural way. One should note that the left and the right R-module actions on $\D_X$ are NOT the same.
\end{remark}

Write $\Mod(\D_X^{\mathrm{op}})$ for the category of right $\D_X$-modules and $\Mod(\D_X)$ for the
category of left $\D_X$-modules.

\begin{theorem}\label{thm:left-right-equivalence}
Tensoring with the canonical bundle yields an equivalence of categories
\[
  \Mod(\D_X) \xrightarrow{\quad\sim\quad} \Mod(\D_X^{\mathrm{op}}),\qquad M\longmapsto \omega_X\otimes_R M,
\]
with quasi-inverse $N\longmapsto \omega_X^{\!*}\otimes_R N$.
\end{theorem}
\begin{proof}
    % 第一部分
    \noindent 1. Let $M \in \Mod(\Dx)$.
    \begin{align*}
        \ox \otimes_R \Dx \otimes_R \oxinv \cdot (\ox \otimes_R M) 
        &= \ox \otimes_R (\Dx \cdot M) \\
        &\subseteq \ox \otimes_R M
    \end{align*}
    \[
        \implies \ox \otimes_R M \in \Mod(\Dx^{\mathrm{op}})
    \]

    % 第二部分
    \noindent 2. Let $N \in \Mod(\Dx^{\mathrm{op}}) = \Mod(\ox \otimes_R \Dx \otimes_R \oxstar)$.
    \begin{align*}
        \Dx \cdot (\oxstar \otimes_R N) 
        &= (\oxstar \otimes_R \ox \otimes_R \Dx \otimes_R \oxstar) \cdot N \\
        &= \oxstar \otimes_R (\ox \otimes_R \Dx \otimes_R \oxstar \cdot N) \\
        &\subseteq \oxstar \otimes_R N
    \end{align*}
    \[
        \implies \oxstar \otimes_R N \in \Mod(\Dx) 
    \]
\end{proof}

In particular, we will henceforth only consider left $\D_X$-modules, written simply as $\D_X$-modules.

\begin{remark}
  For a comprehensive treatment of side change, see \cite[Lemma 1.2.7, Proposition 1.2.9, Proposition 1.2.12]{hotta2008d}.
\end{remark}